\begin{helper}
For every fixed set \(I\), there is a finite family of canonical
pre-terminal history cells \(C_i\), with representatives \(\omega_i\) and
recorded coordinate sets \(R_i\), such that the cells cover exactly the
\(\noderef{TwoBiteProjectionSizeEvent}(I,k,\ell_R,\ell_B)\), are pairwise
disjoint, contain their representatives, have the cylinder description
\[
\omega\in C_i
\quad\Longleftrightarrow\quad
\omega.2.2=\omega_i.2.2
\text{ and }\omega\text{ agrees with }\omega_i\text{ on }R_i,
\]
have \(R_i\) equal to
\(\noderef{TwoBitePreTerminalRecordedPairs}(\omega,\varepsilon,I)\) throughout
the cell, are contained in the projection-size event, and replay
\(\noderef{TwoBiteStageRevealedVertex}\) and
\(\noderef{TwoBiteProjectionPairTouched}\) under any second sample with the
same injection and the same recorded truth table.
\end{helper}
\begin{proof}
Use the refined transcript family supplied by
\noderef{FixedSetPreTerminalRecordedPairsReplayEquality}, whose construction
component is the pre-terminal replay transcript family of
\noderef{FixedSetPreTerminalReplayTranscriptConstruction}.  Thus we have a
finite index type, cells \(C_i\), recorded sets \(R_i\), and representatives
\(\omega_i\).  The cited construction gives the covering equivalence for the
projection-size event, pairwise disjointness, representative membership, the
cylinder description by common injection and agreement on \(R_i\), and the
fact that every cell lies in the projection-size event.  These are exactly the
first, second, third, fourth, and sixth conclusions required here.

It remains to justify the exact-recorded and replay-preservation conclusions
for these cells.  Fix \(i\) and \(\omega\in C_i\).  By the cylinder
description, \(\omega\) has the same injection as \(\omega_i\) and agrees with
\(\omega_i\) on \(R_i\).  The construction of \(R_i\) is the replay output for
\(\omega_i\).  During the \(F_2\)-pass, every queried same-color coordinate is
already a pre-terminal recorded coordinate by
\noderef{FixedSetPreTerminalF2QueriesRecorded}; incident queries after the
staged-revealed set is known are recorded by the incident branch of
\(\noderef{TwoBitePreTerminalRecordedPairs}\); and the later \(X_x(I)\)-pair
queries are recorded by the second branch of the same definition.  Therefore
the replay output has precisely the two-branch characterization of
\(\noderef{TwoBitePreTerminalRecordedPairs}\).  This is exactly the
representative recorded-set conclusion of
\(\noderef{FixedSetPreTerminalRecordedPairsReplayEquality}\), so it identifies
\(R_i\) with
\(\noderef{TwoBitePreTerminalRecordedPairs}(\omega_i,\varepsilon,I)\).

Now apply \noderef{FixedSetPreTerminalReplayUniqueness} to the pair
\((\omega_i,\omega)\).  The common-injection and recorded-truth-table
hypotheses are exactly the cylinder conditions.  Hence the pre-terminal
recorded set for \(\omega\) is the same set \(R_i\).  This proves
\[
e\in\noderef{TwoBitePreTerminalRecordedPairs}(\omega,\varepsilon,I)
\Longleftrightarrow e\in R_i
\]
for every coordinate \(e\), which is the fifth conclusion.

Finally let \(\omega'\) be any second sample with the same injection as
\(\omega\) and the same truth table on \(R_i\).  Since the previous paragraph
identified \(R_i\) with the pre-terminal recorded set of \(\omega\), the
hypotheses of \noderef{FixedSetPreTerminalStageTouchPreservation} apply.
That helper gives equality of
\(\noderef{TwoBiteStageRevealedVertex}\) for every base vertex and equality
of \(\noderef{TwoBiteProjectionPairTouched}\) for every tagged coordinate,
which is the final replay-preservation conclusion.
\end{proof}
